Equal randomization is the gold standard in confirmatory (phase III) clinical trials: each treatment, also called arm, is allocated to approximately the same number of patients, after which some statistical test is performed to identify promising treatments. Using a Response Adaptive Randomization (RAR) procedure, i.e., randomizing patients to treatments based on the observed outcomes for previous patients, has a great potential to make clinical trials more effective: fewer patients may be needed for inferring treatment effects or patient benefit could be increased. Still, the actual use of RAR in clinical trials has been debated for years, as explained  by \cite{Villar23Review}.

There exists many RAR procedures. The first is attributed to \citep{Thompson33}, and variants of this Bayesian strategy have been explored in clinical trials \citep{Thall07}. In the 1950s multi-armed bandits models were introduced to model adaptive selection between different populations \citep{Robbins52Freq} . Since then, bandit algorithms have become a research field of their own \citep{BanditBook}. However, many of these adaptive algorithms are actually \emph{not} randomized hence do not qualify as RAR procedures. Moreover, the type of performance measures considered in that literature are not the typical testing metrics common in clinical trials. Another class of RAR procedures are urn-based designs, such the Randomized Play the Winner (RPW) \citep{WeiDurham78RPW} that was actually used in a real trial \citep{Bartlett85ECMO} but with mixed results as the resulting allocation was extremely unbalanced. In the 2000s, a family of RAR procedures based on a so-called \emph{optimal allocation proportion} emerged \citep{Rosenberger01Opt}. Given some assumptions on the arms distributions and some statistical test performed at the end of the trial, these strategies first derive the allocation optimizing some criterion for a fixed power of the test. For example, minimizing the total number of patients needed for a fixed power of a Wald test in the case of two-armed binary outcomes leads to the so-called Neyman allocation, but many other target allocations have been considered under different assumptions, see \citep{pin2024response} for a recent survey. These \emph{target allocations} all have in common that they depend on the unknown parameters of the arms distributions, and thus a crucial component of the resulting strategy is a RAR procedure that, given a target allocation, ensures that the empirical allocation proportions converge to that target.

% A survey of \cite{pin2024response} reviews several target allocations proposed in response-adaptive randomization, those can be derived depending on the underlying optimality criterion, such as maximizing power or improving patient benefit. In two-arm trials this includes, for example, the Neyman allocation for power-optimality and failure-minimizing allocations (RSIHR), while for multi-arm settings (\(K > 2\)) examples include the multi-arm power-optimal allocation of \cite{tymofyeyev2007implementing} and contrast-based optimal allocations of \cite{biswas2011multi}.
%
Our paper is about such \emph{adaptive targeting mechanisms}. Many such mechanisms were originally studied for two arms, such as the celebrated Biased Coin Design of \cite{Efron71BCD}, which targets a $(1/2,1/2)$ allocation. The work of \cite{hu2004asymptotic} studies the Doubly adaptive Biased Coin Design for general (smooth) targets allocations and further extend this design to the multi-armed case ($K \geq 2$). As argued in several works (e.g. \cite{HuRosenberger03VariabilityPower}), empirical allocations that are less variable lead to a better power of the statistical test using the adaptively collected data. This motivates the quest for \textit{asymptotic efficiency}, defined as achieving the minimal asymptotic variance for the allocation proportions subject to a given allocation target. For adaptive targeting mechanisms, this notion was formalized by \cite{hu2006asymptotically}, who derived lower bounds for the asymptotic covariance of allocation proportions. %Achieving efficiency together with randomization and general target allocation functions was always problematic.
While DCBD is actually not asymptotically efficient, \cite{hu2009efficient} introduce the Efficient Randomized Adaptive Design (ERADE), in the two arms setting. Their procedure achieves the asymptotic efficiency bound while preserving randomization and allowing for general target allocation functions, including targets derived from optimality criteria. Under mild regularity conditions, ERADE yields strong consistency and asymptotic normality for both allocation proportions and estimators of treatment effects. However, its extension to more than $2$ arms was left as an open question.



\paragraph*{Contributions} In this paper, we propose an extension of ERADE for $K\geq 2$ arms. More precisely, we introduce a family of algorithms keeping the same spirit as ERADE, which we call $\alpha$-Rebalancing Targeting Strategies ($\alpha$RTS). %Rather than proposing a specific allocation algorithm, we introduce general conditions imposed directly on the allocation probabilities that define broad families of response-adaptive designs.
Interestingly, the $\alpha$RTS framework recovers the extension of ERADE recently proposed by \cite{alkhnefr2025efficient} in a parallel work, but also allows for different algorithmic choices which we explore in the paper. In particular, we show that a Tracking strategy proposed in a different context in the multi-armed bandit literature \citep{garivier2016optimal} is also an $\alpha$RTS. We propose a unified asymptotic analysis of $\alpha$RTS, establishing consistency and asymptotic normality for both allocation proportions and treatment effect estimators, and we show that the resulting allocation proportions achieve the asymptotic efficiency lower bound.

Prior work only apply to target allocations that are strictly positive, and do not explicitly address sparse target regimes in which some treatments are asymptotically eliminated. Yet, in multi-armed setting, such sparse targets have been considered \citep{tymofyeyev2007implementing}. To overcome this limitation, we propose a variant of the $\alpha$RTS family that incorporates an adaptive forced exploration mechanism into the allocation rule. Despite this modification, we show that our proposed $\alpha$RTS-FE inherits the same asymptotic properties  under the same structural conditions and its asymptotic analysis can be extended to sparse target allocations. In this setting, the adaptive exploration mechanism guarantees infinite sampling of all treatments, ensuring the validity of estimator asymptotics even when certain treatments are asymptotically eliminated.

On the technical side, we emphasize that we propose a unified analysis of the $\alpha$RTS and $\alpha$RTS-FE families, that identifies some sufficient conditions for a RAR design to be asymptotically efficient (Lemma~\ref{lem:generic-conditions-main}). These conditions generalize some conditions that can be  extracted from the original ERADE analysis for $K=2$ arms, and we show that they apply to both type of strategies by introducing some appropriate notion of hitting time $\ell_{n,k}$.

In addition to our theoretical contributions, we complement the asymptotic analysis with a comprehensive simulation study designed to assess the finite-sample behavior of the proposed targeting strategies in a trial with $K=3$ arms with binary outcomes, where an homogeneity tests is performed. These experiments investigate both regular and sparse allocation regimes, comparing different algorithmic instantiations within the \(\alpha RTS\) family and their forced-exploration counterparts. In particular, we evaluate convergence properties of allocation proportions and plug-in estimates, as well as the impact on hypothesis testing performance and statistical power in multi-arm settings. The simulations highlight subtle yet practically important differences between designs that are asymptotically equivalent, and illustrate the critical role of forced exploration in ensuring reliable adaptation in sparse regimes.


\paragraph*{Outline}
The remainder of the paper is organized as follows. Section~\ref{sec:notations} introduces the model, notation, and standing assumptions. Section~\ref{sec:alpha-rts} presents the $\alpha$-Rebalancing Targeting Strategies and discusses particular instantiations. Section~\ref{sec:proof_vanilla} present our main theoretical results and sketch their proof. Section~\ref{sec:alpha-rts-fe} introduces the forced-exploration variant and shows that the asymptotic properties are preserved, including in sparse target allocation regimes. Finally in Section~\ref{sec:simulation} we propose some experiments for different algorithms of the two proposed families. Technical lemmas and detailed proofs are given in the appendices.

